% !TeX root = ../main.tex
%%% ---------------------------------------------------------------------------
%%% 相关工作
%%%
%%% 综述不是文献罗列。每一段应当围绕**一条技术脉络**，说清楚
%%% 「大家怎么做的 → 共同的局限 → 与本文的关系」，最后落回本文的切入点。
%%% ---------------------------------------------------------------------------
\section{相关工作}
\label{sec:related}

\subsection{多尺度特征表示}
\label{sec:related:multiscale}

图像金字塔是最早的多尺度方案，代价是重复的前向计算。特征金字塔
网络\cite{lin2017fpn}把多尺度构造转移到特征层面，以近乎不变的计算量获得
了多层级表示，此后成为检测器的标准组件。后续工作从连接拓扑
入手做了大量改进，但横向连接的融合权重在空间上仍是均匀的。

\subsection{注意力机制在检测中的应用}
\label{sec:related:attention}

通道注意力与空间注意力已被广泛用于强化判别性特征。在检测任务中，
注意力多被放在骨干网络内部，用于抑制背景响应；将其用于\emph{层级选择}
——即决定某个空间位置应当更多依赖哪一层级的特征——的研究相对较少，
这正是本文的切入点。

\subsection{遥感图像目标检测}
\label{sec:related:rsdet}

遥感场景的特殊性催生了旋转框表示、角度周期性损失以及大幅面图像的
切片推理策略\cite{wang2021remote}。这些工作与本文正交：本文改进的是
特征融合方式，可以直接嵌入到上述检测框架中。
